% supplementary_material_monotile.tex
%
% Build with:  make monotile-paper    (from the repository root)
% or:          python3 paper/make_strings_facts.py --sections monotile
%              && pdflatex supplementary_material_monotile.tex (x2)
%
% Every number quoted here is recomputed from the pyCICY source by
% paper/make_strings_facts.py and written to figures/monotile_facts.tex,
% which is \input{} below. The prose and the code cannot drift apart
% silently.

\documentclass[11pt,a4paper]{article}

\usepackage[T1]{fontenc}
\usepackage[utf8]{inputenc}
\usepackage[margin=1in]{geometry}
\usepackage[protrusion=true,expansion=false]{microtype}
\usepackage{amsmath,amssymb,amsthm}
\usepackage{graphicx}
\usepackage{booktabs}
\usepackage{enumitem}
\usepackage{xcolor}
\usepackage[hidelinks]{hyperref}
\usepackage[nameinlink,capitalize]{cleveref}
\usepackage{framed}
\usepackage{orcidlink}

\graphicspath{{figures/}}

% Numbers computed by paper/make_strings_facts.py --sections monotile.
% Declared as fallbacks so the document typesets before the script has run;
% the \input below overrides every one with the computed value.
\newcommand{\MTHatEll}{?}\newcommand{\MTTurtleEll}{?}
\newcommand{\MTSpectreEll}{?}\newcommand{\MTHatPlusTurtle}{?}
\newcommand{\MTSpectreFixed}{?}\newcommand{\MTHatEllFloat}{?}
\newcommand{\MTCharPoly}{?}\newcommand{\MTPhiFour}{?}
\newcommand{\MTPhiFourFloat}{?}\newcommand{\MTPhiFourIsRoot}{?}
\newcommand{\MTPhiFourQuadratic}{?}\newcommand{\MTSubTrace}{?}
\newcommand{\MTFreqH}{?}\newcommand{\MTFreqT}{?}
\newcommand{\MTFreqP}{?}\newcommand{\MTFreqF}{?}
\newcommand{\MTFreqHRational}{?}\newcommand{\MTFreqSum}{?}
\newcommand{\MTChiralRatio}{?}\newcommand{\MTChiralIsPhiFour}{?}
\newcommand{\MTReflectedFraction}{?}\newcommand{\MTReflectedTimesOnePlus}{?}
\newcommand{\MTAperiodic}{?}\newcommand{\MTWitness}{?}
\newcommand{\MTWitnessIrr}{?}\newcommand{\MTSitesGeneric}{?}
\newcommand{\MTSitesHat}{?}\newcommand{\MTSitesOff}{?}
\newcommand{\MTBondsHex}{?}\newcommand{\MTOffRatio}{?}
\newcommand{\MTHypSig}{?}\newcommand{\MTClosePivot}{?}
\newcommand{\MTClosePivotRatio}{?}\newcommand{\MTLocSize}{?}
\newcommand{\MTKappa}{?}\newcommand{\MTScanIndices}{?}
\newcommand{\MTIdxAtomic}{?}\newcommand{\MTIdxNegOne}{?}
\newcommand{\MTIdxHalf}{?}\newcommand{\MTNumpyAgree}{?}
\IfFileExists{figures/monotile_facts.tex}{\input{figures/monotile_facts.tex}}{}

\newcommand{\QQ}{\mathbb{Q}}
\newcommand{\ZZ}{\mathbb{Z}}

\theoremstyle{definition}
\newtheorem{remark}{Remark}

\title{\vspace{-1cm}\textbf{What Real-Space Geometry Determines,\\
and What Has No Momentum Space:\\[2pt]
\large Exact Methods for the Aperiodic Monotile Family}}
\author{B.S. Hartshorn \orcidlink{0009-0004-2853-655X} \small \url{https://github.com/brentharts/CICY}}

\begin{document}
\maketitle

\begin{abstract}
Three companion papers asked what a configuration matrix
determines~\cite{SuppMath}, what anomaly cancellation
determines~\cite{SuppStrings}, and what holomorphy
determines~\cite{SuppTwistor}, each drawing a sharp line between the exact
and the unavailable. This paper asks the question of the Hat --- the single
13-sided shape that tiles the plane only aperiodically --- and of the
continuous family Tile$(a,b)$ it belongs to, whose shape parameter recent
work uses as a control knob for the quantum geometric tensor of electrons
placed on its vertices. The answer has an unusual structure: the system's
most famous property is the \emph{absence} of something. An aperiodic tiling
has no Brillouin zone, so every band-structure method is not approximate but
inapplicable, and the topology must be computed in real space. The
real-space invariant of choice, the spectral localizer index, is half the
signature of a matrix --- and a signature is a count of pivot signs,
decidable exactly over any ordered field. We show that the entire chain from
tile geometry to topological index lives in $\QQ(\sqrt3)$, and the
substitution combinatorics in $\QQ(\sqrt5)$, so that every quantity in this
paper is an exact field element or an integer. The inflation factor is
derived as $\varphi^4$, an exact root of an exactly computed characteristic
polynomial; the ratio of unreflected to reflected hats is derived as
$\varphi^4 : 1$; aperiodicity itself is derived from the irrationality of an
eigenvector; the substrate lattice detects the Hat's proportions by an exact
coincidence of coordinates; and the localizer index of a finite patch is
computed by $LDL^{T}$ elimination with no eigenvalue ever taken, then checked
against an independent floating-point diagonalization of the same matrix.
The boundary acquires its fourth formulation: beside what is exact, what is
quoted, and what is not implemented sits what has no momentum space --- a
category in which the standard tool's absence is a property of the system,
and the exact method is not a luxury but the only game in town.
\end{abstract}

\tableofcontents

% =====================================================================

\section{Introduction}
\label{sec:intro}

The series this paper concludes is organized around a single question asked
of four different structures: what determines what, and exactly? For
heterotic compactifications the determining structure was the configuration
matrix~\cite{SuppMath}; for F-theory and orientifolds it was anomaly
cancellation~\cite{SuppStrings}; for tree amplitudes it was
holomorphy~\cite{SuppTwistor}. In each case the useful output was not only
the list of exact quantities but the taxonomy of the rest --- the
distinguishable ways in which a quantity can fail to be available.

The Hat changes the character of the question. In 2023 Smith, Myers, Kaplan
and Goodman-Strauss found the first aperiodic monotile~\cite{Smith2023a}: a
single shape that tiles the plane, but only aperiodically. It sits in a
continuous family Tile$(a,b)$, parametrised by two edge lengths through
$\ell = a/(a+b)$, whose other named members --- the Chevron, the Spectre,
the Turtle, the Comet --- include, at the mirror-symmetric point, the
strictly chiral monotile of the follow-up paper~\cite{Smith2023b}. Recent
work places tight-binding electrons on the vertices of these tilings and
finds graphene-like spectral features, chiral response, flux-periodic
Hofstadter spectra~\cite{Schirmann2024}, and --- the observation this paper
builds on --- that the shape parameter $\ell$ acts as a real-space control
knob for the quantum geometric tensor, driving topological transitions with
no momentum space anywhere in the problem~\cite{RocheCarrasco2025}.

That last clause is the invitation. A crystal's topological invariants are
integrals over the Brillouin zone; an aperiodic tiling has none, so the
invariants must be computed in real space, and the modern real-space
invariant --- the spectral localizer index of Loring and
Schulz-Baldes~\cite{Loring2015, LoringSchulzBaldes2017} --- is \emph{half the signature
of a matrix}. A signature is a count of pivot signs in an $LDL^T$
decomposition. It is decidable, exactly, over any ordered field, with no
eigenvalue ever taken. The question for a package built on exact arithmetic
is then whether the whole chain
%
\begin{equation}
\text{tile geometry} \;\to\; \text{vertex coordinates} \;\to\;
\text{Hamiltonian} \;\to\; \text{localizer} \;\to\; \text{index}
\label{eq:chain}
\end{equation}
%
can be kept inside such a field. It can. The tiles live on a lattice whose
bond directions are multiples of $30^\circ$, so every coordinate and every
hopping amplitude is $p + q\sqrt3$ with $p, q \in \QQ$; the substitution
combinatorics lives in $\QQ(\sqrt5)$, where the golden ratio does; and both
fields are ordered with a decidable order, which is the one property a
signature computation needs.

\Cref{sec:family} sets out the family and its mirror.
\Cref{sec:inflation,sec:chirality} derive the inflation factor, the tile
frequencies, the chirality imbalance and aperiodicity itself from one
quoted matrix. \Cref{sec:substrate} shows the substrate lattice detecting
the Hat's proportions by an exact coincidence.
\Cref{sec:signature,sec:localizer} build the exact signature machinery and
apply it to the localizer, with an independent floating-point route as the
cross-check. \Cref{sec:boundary} restates the boundary in its fourth form.

% =====================================================================

\section{The family, by parameter}
\label{sec:family}

Tile$(a,b)$ depends on its two edge lengths only through
$\ell = a/(a+b) \in [0,1]$, and the named members sit at exact values of it:
the Chevron at $0$, the Hat at $1/(1+\sqrt3)$, the Spectre at $1/2$, the
Turtle at $\sqrt3/(1+\sqrt3)$, the Comet at $1$. Rationalising,
%
\begin{equation}
\ell_{\mathrm{Hat}} = \MTHatEll \approx \MTHatEllFloat ,
\qquad
\ell_{\mathrm{Turtle}} = \MTTurtleEll ,
\qquad
\ell_{\mathrm{Hat}} + \ell_{\mathrm{Turtle}} = \MTHatPlusTurtle .
\label{eq:ells}
\end{equation}

Exchanging $a \leftrightarrow b$ reflects the tile, so
$\ell \mapsto 1 - \ell$ is a mirror on the family. \Cref{eq:ells} then says
the Hat and the Turtle are mirror partners, and the Spectre, at
$\ell = \MTSpectreEll$, is the fixed point (\MTSpectreFixed). That is the
parameter-space shadow of the Spectre's special role: it is the member that
tiles using a single handedness --- the ``vampire einstein'' that casts no
reflection --- while the Hat, as \cref{sec:chirality} quantifies, needs its
mirror image. The mirror map slots into the same catalogue of
involutions the earlier papers used: the orientifold involution
of~\cite{SuppStrings} and the parity conjugation of~\cite{SuppTwistor} are
the same shape of structure, an exact involution on a parameter space with
the physically distinguished object at or near its fixed locus.

% =====================================================================

\section{The inflation factor, derived}
\label{sec:inflation}

The Hat tiling is generated by four metatiles $H, T, P, F$ and a
substitution rule, and the substitution matrix
%
\begin{equation}
M \;=\;
\begin{pmatrix}
3 & 1 & 2 & 2\\
1 & 0 & 0 & 0\\
3 & 0 & 1 & 1\\
3 & 0 & 2 & 3
\end{pmatrix}
\label{eq:subst}
\end{equation}
%
is one of the two quoted inputs in this paper, taken
from~\cite{Smith2023a}. Everything after it is derived, and the two
published numbers it must reproduce are what the derivation is tested
against.

The characteristic polynomial of \cref{eq:subst}, computed over $\ZZ$ by
exact Leverrier--Faddeev with nothing numerical anywhere, has coefficients
$(\MTCharPoly)$, which factor as
%
\begin{equation}
\chi_M(x) \;=\; (x^2 - 7x + 1)(x^2 - 1) .
\label{eq:charpoly}
\end{equation}
%
The candidate Perron eigenvalue $\varphi^4 = \MTPhiFour$, with $\varphi$ the
golden ratio, is then \emph{verified} as a root by exact Horner evaluation
in $\QQ(\sqrt5)$ (\MTPhiFourIsRoot) and satisfies the quadratic factor
$x^2 = 7x - 1$ (\MTPhiFourQuadratic); no eigenvalue solver is involved. One
inflation step multiplies areas by
$\varphi^4 \approx \MTPhiFourFloat$ --- an exact element of $\QQ(\sqrt5)$,
not a fitted number. Its Galois conjugate $\varphi^{-4}$ is the other root
of the same quadratic, their product is the constant term $1$, and the trace
$\MTSubTrace$ of \cref{eq:subst} accounts for all four roots as
$\varphi^4 + \varphi^{-4} + 1 - 1$.

% =====================================================================

\section{Chirality and aperiodicity, from one eigenvector}
\label{sec:chirality}

The Perron eigenvector of \cref{eq:subst} gives the relative frequencies of
the metatiles, computed here by exact elimination in $\QQ(\sqrt5)$:
%
\begin{equation}
f_H = \MTFreqH, \qquad
f_T = \MTFreqT, \qquad
f_P = \MTFreqP, \qquad
f_F = \MTFreqF ,
\label{eq:freqs}
\end{equation}
%
summing to $\MTFreqSum$ and satisfying $M f = \varphi^4 f$ component by
component, which the tests check directly against \cref{eq:subst}. Two
things about \cref{eq:freqs} deserve notice. The $H$ frequency is
\emph{rational} --- exactly one third (\MTFreqHRational) --- a curiosity
worth pinning because it means the aperiodicity argument below cannot rest
on it. And the $T$ frequency is not, which is where that argument does
rest.

\subsection{The chirality imbalance}

The second quoted input is the hat content of each metatile: $H$ contains
four hats of which one is reflected, and $T$, $P$, $F$ contain one, two and
two, all unreflected~\cite{Smith2023a}. Every anti-hat in the tiling
therefore lives inside an $H$ supertile, and weighting \cref{eq:freqs} by
these counts gives the ratio of unreflected to reflected hats:
%
\begin{equation}
\frac{4f_H + f_T + 2f_P + 2f_F - f_H}{f_H}
\;=\; \MTChiralRatio \;=\; \varphi^4 : 1
\qquad (\text{exactly: } \MTChiralIsPhiFour).
\label{eq:chirality}
\end{equation}
%
The Hat tiling is \emph{almost} chiral: one reflected tile per
$\varphi^4 \approx \MTPhiFourFloat$ unreflected ones, and the imbalance is
not a decimal but an element of $\QQ(\sqrt5)$. The reflected fraction of all
hats is $\MTReflectedFraction$, which multiplied by $1 + \varphi^4$ gives
$\MTReflectedTimesOnePlus$ --- i.e.\ it is exactly $1/(1+\varphi^4)$,
rationalised. The published value of \cref{eq:chirality} is reached here
without being quoted, which is the same discipline the earlier papers
applied to the non-Higgsable clusters and the Schubert cell counts: the
tabulated number is the test, not the input. At the other end of the mirror
of \cref{sec:family}, the Spectre needs no reflected tiles at all, and the
family interpolates --- which is precisely the mechanism
of~\cite{RocheCarrasco2025} for switching off the anti-hat-bound zero modes
by moving $\ell$.

\subsection{Aperiodicity as irrationality}

A tiling with a period is a periodic arrangement of a fundamental domain; in
a fundamental domain every metatile occurs an integer number of times; so in
a periodic tiling every relative frequency is rational. The $\MTWitness$
frequency in \cref{eq:freqs} has irrational part $\MTWitnessIrr \sqrt5$, so
no period exists (\MTAperiodic). The headline fact about the Hat is, at this
level of description, \emph{the irrationality of an eigenvector} --- the
same shape of argument that quasicrystal diffraction rests on, and a
pleasingly small amount of machinery for a statement that made newspapers.
What this derivation does not do is prove the Hat tiles the plane at all;
that is the content of~\cite{Smith2023a} and enters through
\cref{eq:subst}. What it derives is that \emph{given} the substitution
structure, periodicity is impossible.

% =====================================================================

\section{The substrate detects the Hat}
\label{sec:substrate}

Every Tile$(a,b)$ is a union of eight kites whose two bond lengths are the
parameters themselves: hexagon centre to edge midpoint ($a$) and edge
midpoint to hexagon vertex ($b$), meeting at right angles along directions
that are multiples of $30^\circ$. The cosines of those directions lie in
$\tfrac12\ZZ[\sqrt3]$, so every vertex coordinate of every patch is exact.

One hexagon's worth of kites, at generic proportions, has
$\MTSitesGeneric$ distinct sites: one centre, six edge midpoints, and
\emph{twelve} outer corners, because the corners contributed by neighbouring
kites do not coincide. At $a : b = \sqrt3 : 1$ --- the Hat's proportions,
by \cref{eq:ells}, and only there --- the twelve merge in pairs into the six
vertices of the $[3.4.6.4]$ Laves tiling, and the count drops to
$\MTSitesHat$. Whether two exact coordinates coincide is decided by the
order on $\QQ(\sqrt3)$, not by a tolerance: at $a : b = \MTOffRatio$, a
ratio within $2\%$ of $\sqrt3$, the count is $\MTSitesOff$ again. The bond
count is $\MTBondsHex$ in every case, since merging vertices does not merge
bonds. The substrate itself, in other words, detects the Hat point as an
exact coordinate coincidence --- a small structural echo of how the
Weierstrass model of~\cite{SuppStrings} detects its special fibres as exact
vanishing-order coincidences.

The patch so built is periodic, and the module says so. Generating a genuine
aperiodic vertex patch requires implementing the substitution
\emph{geometry} --- the placements, not just the counts, of
\cref{eq:subst} --- which is not done here. What the patch supplies is the
local geometry every Tile$(a,b)$ model inherits, and the testbed on which
the machinery of the next two sections runs end to end.

% =====================================================================

\section{A signature, exactly}
\label{sec:signature}

The invariant of \cref{sec:localizer} is the signature of a symmetric
matrix over $\QQ(\sqrt3)$, so this section is about computing signatures
over an ordered field, which is elementary and worth being careful about
anyway.

$LDL^T$ elimination with symmetric pivoting reduces a symmetric matrix to a
diagonal of pivots whose signs give the signature, by Sylvester's law of
inertia. Over $\QQ(\sqrt d)$ the one non-obvious ingredient is the sign of
$p + q\sqrt d$ itself, which is decided by comparing $p^2$ with $d q^2$ ---
a comparison of rationals. The implementation's test suite includes the
close call $26 - 15\sqrt3$, whose sign ($\MTClosePivot$) is decided by
$\MTClosePivotRatio$: positive by one part in $676$. A floating-point
computation gets this right too, but not \emph{provably}, and the difference
between the two is the difference between a numerical experiment and a
certificate.

The degenerate case is a symmetric block whose diagonal is entirely zero
--- the matrix $\bigl(\begin{smallmatrix}0&1\\1&0\end{smallmatrix}\bigr)$ in
miniature. The first implementation handled it by eliminating around the
off-diagonal block and was wrong, which the hand-checkable test caught; the
replacement adds row $j$ to row $i$ and column $j$ to column $i$, a
congruence $S \mapsto P^T S P$ that Sylvester's law \cite{Sylvester1852} guarantees preserves the
signature and that manufactures the non-zero diagonal entry $2S_{ij}$
directly. The hyperbolic block then reports signature and nullity
$\MTHypSig$, as it must. Recording the failed first attempt is in the
spirit of the series: every error found in four papers' worth of modules was
found by an exact check with a hand-computable answer, and this one was the
monotile module's contribution to the pattern.

% =====================================================================

\section{The localizer index, with no floating point computations}
\label{sec:localizer}

On the patch of \cref{sec:substrate} we place the standard two-orbital
Chern insulator --- the Qi--Wu--Zhang model~\cite{QWZ2006} transplanted to a
graph, the usual move for amorphous and quasicrystalline
topology~\cite{Agarwala2017, Resta2011} --- with on-site term $M\sigma_z$ and, along a
bond of direction $(\cos\theta, \sin\theta)$, the hopping
$\tfrac{t}{2}\bigl[i(\cos\theta\,\sigma_x + \sin\theta\,\sigma_y) -
\sigma_z\bigr]$. On this substrate every $\cos\theta$ and $\sin\theta$ is in
$\tfrac12\ZZ[\sqrt3]$, so with rational $M$ and $t$ the Hamiltonian is a
complex Hermitian matrix whose real and imaginary parts are exact.

The spectral localizer of Loring and
Schulz-Baldes~\cite{LoringSchulzBaldes2017} bundles $H$ with the position
operators,
%
\begin{equation}
L \;=\;
\begin{pmatrix}
H - E_0 & \kappa\,\Pi \\
\kappa\,\Pi^\dagger & -(H - E_0)
\end{pmatrix},
\qquad
\Pi = (X - x_0) - i (Y - y_0),
\label{eq:localizer}
\end{equation}
%
and the local index at $(x_0, y_0, E_0)$ is half its signature. The complex
Hermitian $L$ is doubled to the real symmetric
$\bigl(\begin{smallmatrix}A & -B\\ B & A\end{smallmatrix}\bigr)$, which doubles the
signature again, so the index is a quarter of what the elimination of
\cref{sec:signature} returns. On the one-hexagon patch the doubled localizer
is a $\MTLocSize \times \MTLocSize$ matrix over $\QQ(\sqrt3)$, and with
$\kappa = \MTKappa$ the mass scan gives
%
\begin{center}
\begin{tabular}{lcccc}
\toprule
$M$ & $-4$ & $-1$ & $1/2$ & $4$ \\
\midrule
index & $\MTIdxAtomic$ & $\MTIdxNegOne$ & $\MTIdxHalf$ & $\MTIdxAtomic$ \\
\bottomrule
\end{tabular}
\end{center}
%
trivial in both atomic limits and carrying Chern number $\mp 1$ in between,
with every zero-pivot count zero, so the localizer gap is open at every
sampled point and the index is well defined there. Each jump in the row
brackets an exact topological transition between rational masses ---
located without a band structure, because there is none to compute.

The cross-check is the point of the exercise, and it is the sharpest
available: the same matrix is handed to LAPACK, which diagonalises it in
floating point, and the eigenvalue sign count is compared with the exact
signature. One route is field arithmetic and pivot signs; the other is
numerical linear algebra; they share nothing, and they agree at every mass
(\MTNumpyAgree). This is the fourth paper's instance of the series'
methodological refrain --- a quantity computed once is a result, computed
twice by disjoint routes it is a test --- with the new twist that here one
of the routes is exact and the other is the field's standard numerical
practice, so the comparison also calibrates the practice.

% =====================================================================

\section{The boundary, in its fourth form}
\label{sec:boundary}

Each paper in this series has ended by classifying the ways a quantity can
fail to be available, and each system has demanded a new category. Here the
classification reads:

\begin{enumerate}[leftmargin=2em]
\item \emph{Exact and complete.} The parameter values of the named tiles,
the mirror structure, the inflation factor, the metatile frequencies, the
chirality imbalance, the aperiodicity argument, and the localizer index of
any finite patch --- the full chain \cref{eq:chain}, with no tolerance
anywhere in it.

\item \emph{Quoted.} Two combinatorial facts about the published
construction: the substitution matrix \cref{eq:subst} and the hat content
of the metatiles. The derivations downstream reproduce the published
$\varphi^4$ and $\varphi^4 : 1$ without consulting them, which is the test.

\item \emph{Not implemented.} The substitution geometry --- the exact
placements that would generate genuinely aperiodic vertex patches --- and
with it everything requiring large patches or a thermodynamic limit:
spectral functions, the quantum metric maps of~\cite{RocheCarrasco2025},
disorder averages. The finite patches here are cut from the periodic
substrate, and the module says so rather than presenting them as more.

\item \emph{Has no momentum space.} The Brillouin zone, and everything
built on it: Bloch bands, momentum-space Berry curvature, the
$k$-space Chern integral. For an aperiodic tiling these are not
approximations that degrade but tools that do not apply, and the honest
response is not to regret them but to notice that the real-space invariant
was never second best --- it is the definition that survives, and it is
\emph{more} exactly computable than its momentum-space counterpart, being a
signature rather than an integral.
\end{enumerate}

\noindent
The fourth category completes a progression. In~\cite{SuppMath} the
boundary was of the data; in~\cite{SuppStrings} it separated what
consistency forces from what a choice leaves open, and contained a quantity
that \emph{does not exist}; in~\cite{SuppTwistor} the boundary was of the
method, with a relation exact but invisible to a rank. Here the absent
object is the standard method's entire arena, and the system is the better
for it: the constraint that forced the computation into real space is the
same constraint that made it exactly decidable.

% =====================================================================

\section{Conclusion}
\label{sec:conclusion}

The first paper found that a configuration matrix determines more than one
expects; the second, that consistency does; the third, that holomorphy
does. This one finds that real-space geometry does --- and that when the
geometry is aperiodic, real space is not a fallback but the native habitat
of the invariants, where they are signatures of matrices over ordered
fields and therefore decidable.

The specific results are small and sharp. The inflation factor of the Hat
substitution is $\varphi^4$, as an exact root of an exactly computed
polynomial. The tiling's chirality imbalance is $\varphi^4 : 1$, derived
from an eigenvector over $\QQ(\sqrt5)$. Its aperiodicity is the
irrationality of that eigenvector. The Laves substrate detects the Hat's
proportions as an exact coincidence of coordinates, $\MTSitesGeneric$
sites collapsing to $\MTSitesHat$. And the topological index of a finite
patch is a quarter of an exact matrix signature, computed with no floating
point and confirmed by an independent LAPACK diagonalization at every
sampled mass.

The refrain of the series closes on its strongest instance. A quantity
computed once is a result; computed twice, by routes sharing nothing, it is
a test. Here the two routes are an exact elimination and the field's
standard numerical practice, so the disagreement that never came would have
indicted one of them --- and the agreement that did come is a certificate
for both.

% =====================================================================

\section{Source code and limitations}
\label{sec:source}
\footnotesize
Everything above is implemented in \texttt{pyCICY.monotile}, with
\texttt{tests/test\_monotile.py} executing the checks quoted here and
\texttt{examples/monotile.py} reproducing the tables. The numbers in this
document are emitted by \texttt{paper/make\_strings\_facts.py} with
\texttt{--sections monotile}. The repository is at
\url{https://github.com/brentharts/CICY}.

Limitations:
%
\begin{itemize}[leftmargin=2em,itemsep=1pt]
\item The substitution matrix and the hats-per-metatile counts are quoted;
the tiling property itself (that the Hat tiles the plane, and only
aperiodically) is the theorem of~\cite{Smith2023a} and is not re-proved.
\item Finite patches are cut from the periodic Laves substrate; the
substitution geometry that would generate aperiodic patches is not
implemented yet.
\item The localizer is evaluated at exact sample points; no attempt is made yet
to certify the index over intervals of $M$, which would require tracking
the localizer gap as a function of the mass.
\item One tile family, one model: Tile$(a,b)$ with a Qi--Wu--Zhang
Hamiltonian. Hofstadter physics on the tiling, the quantum metric, and the
topological Anderson transitions of~\cite{RocheCarrasco2025} are absent.
\item Exactness has a cost: the $LDL^T$ elimination over $\QQ(\sqrt3)$ on
the $\MTLocSize$-dimensional localizer takes seconds where LAPACK takes
milliseconds. The claim is decidability, not speed.
\end{itemize}

% =====================================================================

\begin{thebibliography}{99}
\setlength{\itemsep}{1pt plus 0.5pt}

\bibitem{SuppMath}
B.S.~Hartshorn,
\emph{What a configuration matrix determines, and what it does not: exact
methods for heterotic compactifications on CICYs},
Pre-print (2026), \url{https://doi.org/10.5281/zenodo.21844007}.

\bibitem{SuppStrings}
B.S.~Hartshorn,
\emph{What anomaly cancellation determines, and what does not exist: exact
methods for F-theory and Type IIB orientifolds},
Pre-print (2026), \url{https://doi.org/10.5281/zenodo.21865768}.

\bibitem{SuppTwistor}
B.S.~Hartshorn,
\emph{What holomorphy determines, and what a rank cannot see: exact methods
for twistor geometry and tree amplitudes},
Pre-print (2026), \url{https://doi.org/10.5281/zenodo.21892738}.


\bibitem{Smith2023a}
D.~Smith, J.~S.~Myers, C.~S.~Kaplan and C.~Goodman-Strauss,
\emph{An aperiodic monotile},
Combinatorial Theory \textbf{4} (2024) 6,
\href{https://arxiv.org/abs/2303.10798}{arXiv:2303.10798}.

\bibitem{Smith2023b}
D.~Smith, J.~S.~Myers, C.~S.~Kaplan and C.~Goodman-Strauss,
\emph{A chiral aperiodic monotile},
Combinatorial Theory \textbf{4} (2024) 13,
\href{https://arxiv.org/abs/2305.17743}{arXiv:2305.17743}.

\bibitem{Schirmann2024}
J.~Schirmann, S.~Franca, F.~Flicker and A.~G.~Grushin,
\emph{Physical properties of an aperiodic monotile with graphene-like
features, chirality, and zero modes},
Phys.\ Rev.\ Lett.\ \textbf{132} (2024) 086402,
\href{https://arxiv.org/abs/2307.11054}{arXiv:2307.11054}.

\bibitem{RocheCarrasco2025}
H.~Roche Carrasco, J.~Schirmann, A.~Mordret and A.~G.~Grushin,
\emph{Family of aperiodic tilings with tunable quantum geometric tensor},
\href{https://arxiv.org/abs/2505.13304}{arXiv:2505.13304}.

\bibitem{LoringSchulzBaldes2017}
T.~A.~Loring and H.~Schulz-Baldes,
\emph{Finite volume calculation of K-theory invariants},
New York J.\ Math.\ \textbf{23} (2017) 1111,
\href{https://arxiv.org/abs/1701.07455}{arXiv:1701.07455}.

\bibitem{Loring2015}
T.~A.~Loring,
\emph{K-theory and pseudospectra for topological insulators},
Ann.\ Phys.\ \textbf{356} (2015) 383,
\href{https://arxiv.org/abs/1502.03498}{arXiv:1502.03498}.

\bibitem{QWZ2006}
X.-L.~Qi, Y.-S.~Wu and S.-C.~Zhang,
\emph{Topological quantization of the spin Hall effect in two-dimensional
paramagnetic semiconductors},
Phys.\ Rev.\ B \textbf{74} (2006) 085308,
\href{https://arxiv.org/abs/cond-mat/0505308}{arXiv:cond-mat/0505308}.

\bibitem{Agarwala2017}
A.~Agarwala and V.~B.~Shenoy,
\emph{Topological insulators in amorphous systems},
Phys.\ Rev.\ Lett.\ \textbf{118} (2017) 236402,
\href{https://arxiv.org/abs/1607.05457}{arXiv:1607.05457}.


\bibitem{Resta2011}
R.~Bianco and R.~Resta,
\emph{Mapping topological order in coordinate space},
Phys.\ Rev.\ B \textbf{84} (2011) 241106,
\href{https://arxiv.org/abs/1111.5697}{arXiv:1111.5697}.

\bibitem{Sylvester1852}
J.~J.~Sylvester,
\emph{A demonstration of the theorem that every homogeneous quadratic
polynomial is reducible by real orthogonal substitutions to the form of a
sum of positive and negative squares},
Philos.\ Mag.\ \textbf{4} (1852) 138.
\url{https://doi.org/10.1080/14786445208647087}

\end{thebibliography}

\end{document}
